% ============================================================================
%  Chapter 4. Rule-Based Candidate Generation and the Exemplar Selection Failure
%
%  \input file. Requires booktabs, amsmath, natbib in the parent preamble.
%  Bibliography: accord/thesis/refs.bib
%
%  Traceability: accord/phase3_mining.py (rules, exemplar measurement, splits),
%  accord/dictionaries.py (drug classes, dose parsing), accord/CORRECTIONS.md
%  (C3, C6, C8).
% ============================================================================

\chapter{Rule-Based Candidate Generation and the Exemplar Selection Failure}
\label{ch:candidates}

\section{The Problem This Chapter Solves, and the One It Creates}
\label{sec:cand-overview}

True medication contradictions between a nursing note and a discharge summary
are rare. Chapter~\ref{ch:data} produced $9{,}783$ admissions; a board-certified
pharmacist, annotating a stratified sample drawn from the candidate pool, found
$12$ genuine contradictions in $299$ pairs. Exhaustive human review of all
$9{,}783$ admissions is not feasible, and exhaustive machine scoring of all
cross-note medication pairs is both computationally wasteful and statistically
hostile: the positive class would sit somewhere below one in ten thousand.

Candidate generation is the standard response. A cheap, high-recall,
low-precision rule set narrows the space to a pool that a model or a human can
examine. The rules are not intended to be accurate. They are intended to be
\emph{transparent} and \emph{enriching}: transparent because each candidate
carries a reason, and enriching because the base rate inside the pool should be
far above the base rate outside it.

This chapter constructs four such rules, reports their yield, and then reports
the finding that dominates the chapter and, arguably, the thesis. \emph{The
mechanism that decided which sentence pair to attach to a candidate was
independent of the mechanism that decided the candidate was a conflict.} A
candidate was flagged because the \emph{set} of actions observed for a drug in
one note conflicted with the set observed in the other, but the sentence pair
written into the candidate record was the first mention of that drug on each
side, which need not be, and $41.4\%$ of the time was not, the mention carrying
the conflicting action.

Every consumer of the candidate pool downstream saw only the sentence pair. The
pharmacist annotated the sentence pair. The weak labeller of
Chapter~\ref{ch:weak} read the sentence pair. The model of
Chapter~\ref{ch:arch} was trained on the sentence pair. For a large minority of
candidates, all three were reasoning about evidence that did not contain the
conflict that had caused the candidate to exist.

\section{Problem Formulation}
\label{sec:cand-formulation}

Let an admission $a$ have a nursing note $N_a$ and a discharge summary $D_a$.
Chapter~\ref{ch:data}'s extractor yields, for each note, a set of medication
mentions. Write $M(N_a)$ and $M(D_a)$ for these sets. Each mention
$m$ carries a normalised drug $\mathrm{drug}(m)$, an action label
$\mathrm{act}(m) \in \{\textsc{cess}, \textsc{cont}, \textsc{init},
\textsc{dose}, \textsc{unc}\}$, an optional numeric dose
$(\mathrm{val}(m), \mathrm{unit}(m))$, a sentence $\mathrm{sent}(m)$, and, in
the rebuilt pipeline only, a character offset $\mathrm{pos}(m)$ into the note.

A \emph{candidate} is a tuple
\[
c = \bigl(a,\; r,\; d,\; m_N,\; m_D\bigr),
\]
where $r \in \{1,2,3,4\}$ is the rule that produced it, $d$ is the drug or drug
pair at issue, and $(m_N, m_D)$ is the \emph{exemplar pair}: the single mention
from each note that is presented as evidence.

The distinction that this chapter turns on is between the \emph{trigger
predicate} $T_r(a, d)$, which decides whether a candidate exists, and the
\emph{exemplar selector} $E_r(a, d)$, which decides what $(m_N, m_D)$ is. A
correct implementation requires
\[
T_r \text{ holds} \;\Longrightarrow\; \mathrm{act}(m_N) \text{ and }
\mathrm{act}(m_D) \text{ are the two actions that made } T_r \text{ hold}.
\]
The reference implementation does not satisfy this, because $T_r$ quantifies
over sets of mentions and $E_r$ was implemented as
\texttt{.iloc[0]}, the first row of the drug's mention group.

\section{The Four Rules}
\label{sec:rules}

\subsection{Rule 1: direct action conflict}

Rule 1 fires when the same drug $d$ appears in both notes and the action sets
disagree in a specific direction:
\[
T_1(a,d) \;=\;
\bigl(\textsc{cont} \in A_N \wedge \textsc{cess} \in A_D\bigr)
\;\vee\;
\bigl(\textsc{cess} \in A_N \wedge \textsc{cont} \in A_D\bigr),
\]
where $A_N = \{\mathrm{act}(m) : m \in M(N_a), \mathrm{drug}(m) = d\}$ and $A_D$
is defined analogously. This is the canonical case: the nursing note records the
drug as held and the discharge summary orders it continued, or the reverse.

Note the existential quantifier. Rule 1 fires if \emph{some} mention on each
side carries the required action. A drug mentioned five times in a discharge
summary with four continuation mentions and one cessation mention satisfies the
predicate. This is deliberate and correct for a high-recall trigger. It is also
precisely what makes the exemplar selection problem non-trivial.

\subsection{Rule 2: drug class conflict}

Rule 2 generalises Rule 1 to the pharmacologic class. Two \emph{different} drugs
in the same class, one continued and one ceased across the two notes, may
represent a substitution that was documented on one side and not the other. It
uses a hand-written dictionary of $10$ classes covering $31$ drugs
(ACE inhibitors, ARBs, beta blockers, loop diuretics, anticoagulants, statins,
insulins, nitrates, antiarrhythmics, aldosterone antagonists). The predicate is
$T_1$ with the drug identity relaxed to class identity, plus a guard that the
drug sets on the two sides are not identical, which would return the candidate
to Rule 1's territory.

The class dictionary is a limitation in itself. Ten classes and $34$ drugs is a
small fraction of the pharmacologic space, it is not derived from a standard
terminology such as ATC or RxNorm, and its class boundaries encode one author's
judgement. It has no mechanism for cross-class interactions, which are a large
share of clinically significant medication problems.

\subsection{Rule 3: dosage conflict}

Rule 3 fires when the same drug carries a parseable numeric dose in both notes,
the two doses reduce to the same canonical unit, and they differ by more than
$20\%$ of the larger value:
\[
T_3(a,d) \;=\; \exists\, (m_N, m_D) :\;
\mathrm{unit}(m_N) = \mathrm{unit}(m_D)
\;\wedge\;
\frac{\max - \min}{\max} > 0.20 .
\]
The $20\%$ threshold is a judgement call with no empirical derivation. It is not
tuned, and it is not clinically grounded in the sense that different drug
classes tolerate very different dose variation. A $25\%$ change in a statin dose
and a $25\%$ change in a warfarin dose are not comparable events. This is a
known weakness of the rule and it is not addressed anywhere in this thesis.

Rule 3's exemplar selection is correct by construction: the exemplar pair is the
$(m_N, m_D)$ pair achieving the largest percentage difference, so it necessarily
carries the conflict. There is no \texttt{.iloc[0]} defect to fix here, because
the rule had no reference implementation at all. It exists only in the rebuild,
for the reason given in Chapter~\ref{ch:data}: the reference extractor stored
dose strings and never parsed them to numbers, so there was nothing to compute a
percentage from.

\subsection{Rule 4: omission conflict}

Rule 4 fires when a drug receives a cessation action in the nursing note and is
\emph{absent} from every extracted mention in the discharge summary, classified
or not:
\[
T_4(a,d) \;=\;
\bigl(\exists\, m \in M(N_a) : \mathrm{drug}(m) = d \wedge
\mathrm{act}(m) = \textsc{cess}\bigr)
\;\wedge\;
d \notin \{\mathrm{drug}(m) : m \in M(D_a)\}.
\]
This is the clinically most interesting rule and the one with the weakest
evidence, because absence is not the same as a decision. A drug held at the
bedside and never mentioned again in the discharge summary may have been
deliberately discontinued and correctly omitted, or may have been dropped from
the discharge list by oversight. The rule cannot distinguish these, and neither
can the sentence pair it produces, because there is no discharge sentence: the
discharge side of a Rule 4 candidate is the literal string \texttt{absent}.

That property is worth flagging now and is returned to in
Chapter~\ref{ch:weak}: Rule 4 is the only rule whose candidates present the
downstream consumer with a structurally different input, a nursing sentence
paired with an empty string. It is also the only rule on which the weak
labeller and the trained model behave differently from the other three.

\section{Candidate Yield}
\label{sec:cand-yield}

\begin{table}[t]
\centering
\small
\begin{tabular}{clrr}
\toprule
\textbf{Rule} & \textbf{Trigger} & \textbf{Candidates} & \textbf{Admissions} \\
\midrule
1 & Direct action conflict (same drug, opposite actions) & 886   & 701 \\
2 & Drug class conflict (same class, different agents)   & 203   & 195 \\
3 & Dosage conflict ($>20\%$ difference)                 & 521   & n/a \\
4 & Omission (held in nursing, absent from discharge)    & 370   & n/a \\
\midrule
  & \textbf{Total raw}                                   & \textbf{1,980} & n/a \\
  & \textbf{After temporal prefilter}                    & \textbf{1,978} & \textbf{1,203} \\
\bottomrule
\end{tabular}
\caption{Candidate yield by rule. Per-rule admission counts were preserved in
the run log only for Rules 1 and 2. The temporal prefilter removed two
candidates in total; Chapter~\ref{ch:data}, Section~\ref{sec:charttime},
explains why this filter should be understood as an instrument failure rather
than a measurement, and why the two-candidate reduction carries no information.}
\label{tab:yield-rules}
\end{table}

Table~\ref{tab:yield-rules} reports the pool. Two observations.

First, the pool is small relative to the cohort: $1{,}978$ candidates from
$9{,}783$ admissions, so roughly one admission in eight contributes a candidate,
and the $1{,}978$ candidates come from only $1{,}203$ distinct admissions, about
$1.64$ candidates per contributing admission. The concentration is not
incidental; Section~\ref{sec:splits} shows that it is the reason a row-level
train/test split leaks.

Second, the drug distribution is dominated by a handful of agents, shown in
Table~\ref{tab:top-drugs}. Furosemide alone accounts for $373$ candidates,
nearly a fifth of the pool.

\begin{table}[t]
\centering
\small
\begin{tabular}{lrl}
\toprule
\textbf{Drug} & \textbf{Candidates} & \textbf{Class} \\
\midrule
Furosemide  & 373 & Loop diuretic \\
Metoprolol  & 297 & Beta blocker \\
Heparin     & 166 & Anticoagulant \\
Warfarin    & 145 & Anticoagulant \\
Amiodarone  & 110 & Antiarrhythmic \\
\bottomrule
\end{tabular}
\caption{The five most frequent drugs in the candidate pool. These five account
for $1{,}091$ of $1{,}978$ candidates. Within Rule 2, anticoagulants contribute
$107$ of $203$ candidates and beta blockers $34$, so two of ten classes supply
$69\%$ of the class-conflict pool. A model trained on this pool sees a small
number of drugs repeatedly and has ample opportunity to key on drug identity
rather than on the relationship between the two sentences.}
\label{tab:top-drugs}
\end{table}

\section{The Exemplar Selection Failure}
\label{sec:exemplar}

\subsection{The defect}

The reference mining code for Rules 1 and 2 has the following shape:

\begin{quote}\ttfamily\small
conflict = ('continuation' in n\_actions and 'cessation' in d\_actions) \\
\hphantom{conflict = }or ('cessation' in n\_actions and 'continuation' in d\_actions) \\
if conflict: \\
\hphantom{if }n\_row = nursing[nursing['drug'] == drug].iloc[0] \\
\hphantom{if }d\_row = discharge[discharge['drug'] == drug].iloc[0]
\end{quote}

The first two lines evaluate the trigger over sets. The last two lines select
the exemplar by positional index into the drug's mention group. Nothing
constrains \texttt{n\_row} and \texttt{d\_row} to carry the actions that made
\texttt{conflict} true.

The failure mode is concrete. Suppose a discharge summary mentions furosemide
five times: four windows resolve to \emph{continuation} and one, later in the
note, resolves to \emph{cessation}. Suppose the nursing note has a single
\emph{continuation} mention. Then $A_N = \{\textsc{cont}\}$,
$A_D = \{\textsc{cont}, \textsc{cess}\}$, the trigger fires on the
$(\textsc{cont}, \textsc{cess})$ branch, and the candidate is real by the rule's
own definition. But \texttt{.iloc[0]} on the discharge side returns a
\emph{continuation} row. The exemplar pair that the pipeline labels a
contradiction consists of two sentences in which the two notes \emph{agree}.

\subsection{Why the defect was invisible}

Chapter~\ref{ch:data} noted that the reference mentions table stored no
positional information. That is the enabling condition. Without a sentence index
or a character offset, mentions of the same drug within the same note are
indistinguishable rows, and \texttt{.iloc[0]} is not obviously different from any
other choice; there was no representation in which ``the mention that carries
the conflict'' could be expressed. The bug is therefore not a typo. It is a
consequence of an upstream schema decision, and it is the clearest instance in
this project of the pattern that \citet{sculley2015debt} describe as
configuration and data-dependency debt: a defect made unreachable by the absence
of the field that would have exposed it.

\subsection{The fix}

The rebuilt implementation (\texttt{accord/phase3\_mining.py}) selects, on each
side independently, among only those mentions whose \emph{own} action equals the
action required by that specific candidate's conflict branch:
\[
E_r^{\text{new}}(a,d) = \bigl(\arg\!\min_{m \in C_N} \delta(m),\;
\arg\!\min_{m \in C_D} \delta(m)\bigr),
\]
\[
C_N = \{m : \mathrm{drug}(m)=d,\ \mathrm{act}(m) = \rho_N\}, \quad
C_D = \{m : \mathrm{drug}(m)=d,\ \mathrm{act}(m) = \rho_D\},
\]
where $(\rho_N, \rho_D)$ is the required action pair for the branch that fired,
and $\delta(m)$ is the character distance from $m$ to the nearest dose-bearing
mention in the same note. The tie-break is a heuristic with a stated rationale:
a mention adjacent to a concrete dose is more likely to be a substantive
instruction than a passing reference in an allergy list or a problem list. If no
dose-bearing mention exists in the note, the tie-break falls back to the
earliest qualifying sentence, for determinism.

The tie-break is not validated. It affects only which of several
\emph{qualifying} mentions is chosen, so it cannot reintroduce the original
defect, but the claim that dose adjacency indicates clinical substance is an
assumption. [TODO: measure whether the dose-adjacency tie-break changes the
pharmacist's judgement relative to an earliest-sentence tie-break, on a sample
of candidates where the two disagree.]

\subsection{The measurement}

The rebuild instruments both selectors and records, for every action-conflict
candidate, whether the old exemplar's action pair matched the required pair and
whether the fix changed the selected rows. Table~\ref{tab:exemplar} reports the
result.

\begin{table}[t]
\centering
\small
\begin{tabular}{lrr}
\toprule
 & \textbf{Old exemplar correct} & \textbf{Exemplar changed by fix} \\
\midrule
Rule 1 (direct action conflict) & 60.5\% & n/a \\
Rule 2 (drug class conflict)    & 50.4\% & n/a \\
\midrule
\textbf{Rules 1 $+$ 2 pooled}   & \textbf{58.6\%} & \textbf{739 / 1,307 (56.5\%)} \\
\bottomrule
\end{tabular}
\caption{Exemplar correctness under the reference selector. ``Old exemplar
correct'' is the fraction of action-conflict candidates for which the
\texttt{.iloc[0]} exemplar pair actually carried the two actions that triggered
the candidate. The complement, $41.4\%$, is the fraction of action-conflict
candidates justified to every downstream consumer by a sentence pair that did
not contain the conflict. Rule 2 is worse than Rule 1, as expected: class-level
conflicts quantify over more mentions, so the probability that the first mention
happens to be the conflicting one is lower. The rebuild mined $1{,}307$
action-conflict candidates against the reference's $1{,}089$
($886 + 203$), because it runs against the corrected cohort of
Chapter~\ref{ch:data} and an offset-bearing mentions table; the rate is measured
on the rebuild's pool and is reported as such.}
\label{tab:exemplar}
\end{table}

Three readings of Table~\ref{tab:exemplar} are worth separating.

\paragraph{As a bug report.} $41.4\%$ of action-conflict candidates carried
evidence inconsistent with their own trigger, and $56.5\%$ of them change
exemplar under the fix. The two numbers differ because a candidate can change
exemplar on one side while the old exemplar happened to be correct on the other,
or vice versa.

\paragraph{As a bound on annotation validity.} The pharmacist annotated the
exemplar pair. For roughly two in five action-conflict candidates, the sentence
pair presented did not contain the flagged conflict. The pharmacist's judgement
on those pairs is a correct judgement about the wrong evidence. This means the
$4.0\%$ contradiction rate reported in Section~\ref{sec:gold} is an
\textbf{upper bound on the validity of the annotation exercise}, not a clean
estimate of the candidate pool's precision. Whether the true precision is higher
or lower is not determinable from the data collected, because we do not know how
the pharmacist would have ruled on the corrected exemplars. This limitation is
restated in every chapter that uses the $4.0\%$ figure.

\paragraph{As a scoping constraint.} The exemplar defect contaminates
\emph{annotation}. It does not contaminate the measurements in
Chapter~\ref{ch:eval} that are computed on model scores against the annotation:
the AUROC, the threshold sweep, the severity correlations, and the transfer
collapse are all statements about whether the model's output ordering carries
information about the pharmacist's label, whatever that label was measured
against. If the model had learned anything from the exemplar sentences, it would
show up as ranking signal against those same sentences. It does not.
Section~\ref{sec:eval-scoping} in Chapter~\ref{ch:eval} makes this argument
formally.

\section{The Gold Standard}
\label{sec:gold}

\subsection{Sampling}

A stratified sample of $300$ candidates was drawn from the $1{,}978$-pair pool,
allocated across rules as shown in Table~\ref{tab:gold}. The allocation
oversamples Rules 2 and 4 relative to their share of the pool, on the reasoning
that these are the rules whose behaviour was least understood.

\subsection{Annotation protocol}

A single board-certified pharmacist annotated $299$ of the $300$ sampled
pairs.\footnote{The sampling frame is $300$; $299$ annotations were returned.
[TODO: identify which stratum lost a row. The per-rule rates in
Table~\ref{tab:gold} are computed against the $300$-row stratum sizes, which
reproduce the reported percentages exactly, so the missing row's stratum cannot
be recovered from the rates alone.]} For each pair the annotator recorded a
three-way label (\emph{contradiction}, \emph{not a contradiction},
\emph{ambiguous}), a free-text justification, a confidence rating, and, for
contradictions, a $1$ to $5$ danger rating.

Annotation was performed in two passes separated by a two-week washout. This
controls for within-session drift and fatigue. It does not establish inter-rater
reliability, because there is one rater.

\subsection{A retracted reliability claim}

The project previously reported a Cohen's $\kappa$ of $1.000$ between ``two
independent pharmacist judgement instruments,'' offered as an answer to reviewer
requests for inter-rater reliability. That claim is withdrawn. The two
``instruments'' were the binary contradiction label and the $1$ to $5$ danger
rating. A pair that is not a contradiction has no danger by definition, so
$\texttt{label} = \textit{contradiction} \iff \texttt{danger} > 0$ is an
identity and $\kappa = 1.000$ is guaranteed by construction
\citep{cohen1960kappa}. It measures nothing about annotator reliability.

The honest statement is the one made here: this study has one annotator, a
two-pass protocol with a two-week washout, and \textbf{no established inter-rater
reliability}. With $12$ positives in $299$ pairs, a second annotator would be
the highest-value single addition to this work, and the absence of one is a
first-order limitation on every conclusion that rests on the gold labels.

\subsection{Results}

\begin{table}[t]
\centering
\small
\begin{tabular}{clrrrr}
\toprule
\textbf{Rule} & \textbf{Trigger} & \textbf{Sampled} & \textbf{True} & \textbf{Rule} & \textbf{ACCORD} \\
              &                  &                  & \textbf{contra.} & \textbf{precision} & \textbf{precision} \\
\midrule
1 & Direct action conflict & 134 & 3 & 2.2\% & 2.3\% \\
2 & Drug class conflict    & 31  & 3 & 9.7\% & 10.0\% \\
3 & Dosage conflict        & 79  & 4 & 5.1\% & 5.2\% \\
4 & Omission               & 56  & 2 & 3.6\% & 2.4\% \\
\midrule
  & \textbf{Pooled}        & \textbf{300} & \textbf{12} & \textbf{4.0\%} & \textbf{4.7\%} \\
\bottomrule
\end{tabular}
\caption{Gold-standard results by rule. ``Rule precision'' is the fraction of
that rule's annotated candidates the pharmacist judged to be true
contradictions, that is, the rule's own precision as a detector. ``ACCORD
precision'' is the precision of the trained detector of Chapter~\ref{ch:arch}
restricted to that rule's candidates, reported here for direct comparison and
derived in full in Chapter~\ref{ch:eval}. The per-rule true-contradiction counts
$(3,3,4,2)$ are derived from the reported rates and stratum sizes and sum to the
$12$ observed positives. Sampled counts are the stratum sizes of the $300$-row
frame; $299$ annotations were returned.}
\label{tab:gold}
\end{table}

Of $299$ annotated pairs: $12$ contradictions ($4.0\%$), $241$ not
contradictions, and $46$ ambiguous ($15.4\%$). The ambiguous rate is itself a
result. One in six candidates could not be adjudicated by a board-certified
pharmacist from the presented evidence. Some of that is the exemplar defect,
some is the inherent underspecification of a two-sentence view of a medication
decision, and the two cannot be separated with the data collected.

\subsection{The finding in Table~\ref{tab:gold}}

Read the last two columns together. For Rules 1, 2, and 3, ACCORD's precision on
a rule's candidates is essentially identical to that rule's own base rate:
$2.3$ against $2.2$, $10.0$ against $9.7$, $5.2$ against $5.1$. For Rule 4 it is
\emph{below} the base rate, $2.4$ against $3.6$.

The interpretation is direct. A detector that adds discriminative value on top
of a rule should have precision above that rule's base rate, because it should
be discarding some of the rule's false positives. ACCORD does not. Its precision
tracks the base rate because it predicts \emph{positive} for essentially every
candidate the rules hand it, in which case precision is arithmetically forced to
approximately the base rate divided by the predicted-positive rate.
Chapter~\ref{ch:weak} shows where that near-universal positive prediction comes
from, and Chapter~\ref{ch:eval} shows that it holds at every threshold, not just
at $0.5$.

The Rule 4 case is the informative exception. Rule 4 is the only rule on which
the model does \emph{not} predict positive almost universally (it predicts
positive at $74.5\%$), and it is the only rule on which the model's precision
falls below the base rate. Arithmetically, ACCORD recovers only one of the two
true Rule 4 contradictions, which accounts for both the precision drop and the
single miss in the $11/12$ recall reported in Chapter~\ref{ch:eval}. Where the
model exercises any selectivity at all, it selects wrongly.

\section{Splitting, and Why the Reference Split Leaked}
\label{sec:splits}

The reference pipeline split the candidate pool at the row level:

\begin{quote}\ttfamily\small
labels\_df = labels[labels['weak\_label'] != 'ambiguous'].sample(frac=1, random\_state=42) \\
train\_df = labels\_df.iloc[:int(n*0.8)]
\end{quote}

With $1{,}978$ candidates drawn from $1{,}203$ admissions, and with a patient
population that has known repeat admissions, this places candidates that share a
nursing note, a discharge summary, an admission, or a patient on both sides of
the split. Two candidates from the same admission differ only in the drug at
issue; the surrounding note text, the de-identification surrogates, and the
documentation style are identical. A model that memorises note-level surface
features will score them together.

The rebuild splits by \texttt{SUBJECT\_ID} with disjointness asserted in code:
\[
S_{\text{train}} \cap S_{\text{val}} = S_{\text{train}} \cap S_{\text{test}}
= S_{\text{val}} \cap S_{\text{test}} = \emptyset,
\qquad
S_{\text{train}} \cup S_{\text{val}} \cup S_{\text{test}} = S .
\]
Patient-level grouping is the correct unit here because it is the only unit that
guarantees no shared note and no shared documentation episode. Assertions rather
than comments are used deliberately: the reference pipeline's split was not
wrong because the author believed row-level splitting was acceptable, it was
wrong because nothing in the code would have complained.

All results reported in Chapters~\ref{ch:arch} and~\ref{ch:eval} for the trained
model use a held-out set that excludes the $299$ annotated pairs. An earlier
reported configuration trained on the full $1{,}978$-pair pool, which contains
those $299$ pairs. That configuration is training on the evaluation set, its
results are not reported anywhere in this thesis, and the reader should treat
any earlier project artifact quoting it as superseded.

\section{Proposed Figures}
\label{sec:cand-figures}

\begin{figure}[t]
\centering
\framebox[\linewidth]{\parbox[c][6cm][c]{0.92\linewidth}{\centering
\textit{[Figure asset 4.1: schematic of the exemplar defect. Two columns,}
\textit{``nursing note'' and ``discharge summary''. Each column shows five}
\textit{abstract mention slots with action tags (no clinical text). The}
\textit{discharge column has four CONT slots and one CESS slot at position 5.}
\textit{An arrow labelled ``trigger: sets conflict'' spans both columns. A}
\textit{second arrow labelled \texttt{.iloc[0]} points at slot 1 (CONT). A}
\textit{third arrow labelled ``fix: action-matched + dose-adjacent'' points at}
\textit{slot 5 (CESS). Caption box: ``the pair shown to the annotator, the LLM,}
\textit{and the model is the \texttt{.iloc[0]} pair''.]}}}
\caption{The exemplar selection defect. The trigger predicate quantifies over
sets of mentions; the exemplar selector indexes positionally. Nothing forces the
two to agree, and for $41.4\%$ of action-conflict candidates they did not. No
clinical text appears in this figure; slots are abstract.}
\label{fig:exemplar}
\end{figure}

\begin{figure}[t]
\centering
\framebox[\linewidth]{\parbox[c][5cm][c]{0.92\linewidth}{\centering
\textit{[Figure asset 4.2: grouped bar chart, one group per rule. Bars: (a)}
\textit{rule precision from the pharmacist, 2.2 / 9.7 / 5.1 / 3.6; (b) ACCORD}
\textit{precision on that rule's candidates, 2.3 / 10.0 / 5.2 / 2.4. A dashed}
\textit{horizontal line at the pooled 4.0\% base rate. Rule 4's pair is the only}
\textit{one where (b) is below (a); highlight it.]}}}
\caption{ACCORD's precision tracks each rule's own base rate. The reader should
notice that the model contributes no lift over the rule that generated the
candidate on any of the four rules, and that on Rule 4, the only rule where it
declines a substantial fraction of candidates, its precision falls below the
base rate.}
\label{fig:rule-precision}
\end{figure}

\section{Limitations}
\label{sec:cand-limitations}

\paragraph{The rules encode one author's clinical judgement.} The class
dictionary, the $20\%$ dose threshold, the choice of four rules rather than
forty, and the decision to treat absence as evidence in Rule 4 are all
unvalidated design choices. No clinician reviewed the rule set before it was
run.

\paragraph{Rules 3 and 4 have no reference implementation.} They exist only in
the rebuild. Any comparison between rebuild and reference behaviour is therefore
restricted to Rules 1 and 2, which is why Table~\ref{tab:exemplar} covers only
those two.

\paragraph{The exemplar defect bounds annotation validity, and the bound is
one-sided in an unknown direction.} We know $41.4\%$ of action-conflict
candidates were annotated against evidence that did not carry their conflict. We
do not know how the pharmacist would have ruled on corrected evidence. Both a
higher and a lower true precision are consistent with the data. Re-annotating a
sample of the changed candidates under the corrected selector is the single
measurement that would resolve this, and it was not performed. [TODO: re-annotate
$n \approx 100$ of the $739$ changed candidates and report the label-change
rate.]

\paragraph{One annotator.} Section~\ref{sec:gold}. Twelve positives is a small
number on which to rest a benchmark, and the confidence intervals in
Chapter~\ref{ch:eval} reflect that.

\paragraph{The candidate pool bounds recall absolutely.} Nothing outside the
$1{,}978$ candidates can ever be detected by this pipeline. The rules' recall
against the full set of true contradictions in the $9{,}783$ admissions is
unmeasured and unmeasurable without exhaustive review. Every recall figure in
this thesis is recall \emph{within} the candidate pool.

\section{Summary}
\label{sec:cand-summary}

Four rules produce $1{,}978$ candidates over $1{,}203$ admissions, enriched to a
pharmacist-measured base rate of $4.0\%$ from a background rate that is orders
of magnitude lower. That much works.

The chapter's principal finding is that the enrichment came with a defect that
was invisible to every downstream consumer. For $41.4\%$ of action-conflict
candidates, the sentence pair carrying the justification did not carry the
conflict, because the trigger quantified over sets while the exemplar selector
indexed positionally, and because the upstream schema had no field in which the
distinction could be expressed. The fix changes the exemplar for $739$ of
$1{,}307$ candidates.

The consequence is a bound, not a correction: the $4.0\%$ figure is an upper
bound on the validity of the annotation exercise rather than an estimate of the
pipeline's precision. Chapter~\ref{ch:eval} shows that this bound does not
compromise the negative result, because the negative result is measured on
model scores against the same annotations, and the model's scores turn out to
carry no information about those annotations at all.

Chapter~\ref{ch:weak} takes the candidate pool and asks a language model to
label it.
